%% notes-tex/010-index-defect/sections/1-what-broke.tex

\section{What broke\secstatus{todo}}
\label{sec:broke}

Experiment 010 trained an equivariant cell graph transformer on 376{,}732 Kuzmin
trigenic records and then scored three unmeasured design spaces with it. The
largest, \texttt{inference\_3}, covered 465{,}735{,}532 triples genome-wide and
is what the panel-12 and panel-24 gene selections were drawn from. Those
selections went to the bench.

Every prediction in that run is wrong, in two ways that compound.

\subsection{The index space is rebuilt, not stored\secstatus{todo}}
\label{sec:rebuilt}

The model addresses genes by position. Its embedding table has 6{,}607 rows, read
directly from the checkpoint, and a record's perturbed genes arrive as integer
positions in the cell graph's sorted node list. Nothing ties the two together.
The cell graph is constructed at run time from whatever gene set the genome
object currently yields, and no assertion compares its node count against the
\texttt{gene\_num} the checkpoint was built for.

So the index space is a derived quantity recomputed at every use, not an artifact
stored alongside the build. When the genome moves, every embedding row silently
comes to mean a different gene.

\subsection{Defect one, a uniform shift of 28\secstatus{todo}}

Between January and February 2026 the genome gene set went from 6{,}607 entries
to 6{,}579. The two \texttt{inference\_2} and \texttt{inference\_3} dataset builds
logged the smaller number at build time. The 28 missing entries are the
mitochondrial open reading frames \texttt{Q0010} through \texttt{Q0297}, and
because they sort before every nuclear systematic name they occupy positions 0
through 27.

The consequence is not partial. Of the 6{,}579 genes present in both sets, not
one keeps its index and every one shifts by exactly 28. \texttt{YAL001C} sits at
row 28 in the space the checkpoint learned and was fed as row 0. Twenty-eight
embedding rows were never addressed at all.

\subsection{Defect two, triples scored as doubles\secstatus{todo}}
\label{sec:doubles}

The triple generator drew its candidate genes from the Costanzo single-mutant
gene list and never intersected them with the model's gene space. Those are
different universes: of the 5{,}493 genes Costanzo measured, 5{,}440 are inside
the model's 6{,}607 and 53 are outside. By our own annotation 35 of the 53 are
absent from the GFF entirely, 11 are typed \texttt{pseudogene}, 6
\texttt{blocked\_reading\_frame} and 1 \texttt{transposable\_element\_gene}.

A gene the index cannot resolve does not raise. The perturbation processor builds
its index list with a set comprehension over node names, so an unmatched name
contributes nothing and the list is simply shorter, and the readout pools by mean
over however many indices arrive. A two-element list is a perfectly valid double.

Among the 432 distinct panel combinations that were stored, 160 had all three
genes indexable, 266 had two and 6 had one. They carry only 220 distinct
prediction values. The collapse is visible without running anything:

\begin{verbatim}
YIL174W   + YKL033W-A + YPL081W  ->  0.711426
YJL017W   + YKL033W-A + YPL081W  ->  0.711426
YKL033W-A + YLR312C-B + YPL081W  ->  0.711426
\end{verbatim}

Three different triples, one value, because all three of \texttt{YIL174W},
\texttt{YJL017W} and \texttt{YLR312C-B} were unresolvable and each triple
collapsed onto the double \texttt{\{YKL033W-A, YPL081W\}}.

\subsection{A name problem, not a missing gene\secstatus{todo}}
\label{sec:alias}

The distinction matters and is easy to state wrongly. Three of the five
unresolvable panel names are simply outdated systematic names:
\texttt{YJL017W} is now \texttt{YJL016W}, \texttt{YKL200C} is now
\texttt{YKL201C}, and \texttt{YLR312C-B} is now \texttt{YLR313C}. R64 records
\texttt{YLR313C} as one verified gene, SGD:S000004305, carrying
\texttt{Alias: ['SPH1', 'YLR312C-B']}.

That gene has a trained embedding row, number 4389, and always did. What the
original run lacked was a lookup entry for the outdated string, because the gene
set it indexed against stores current systematic names. Rescoring under the
resolved name is therefore not inventing a prediction for a gene the model never
saw; it is asking the model about the locus the strain actually deletes.

The remaining two, \texttt{YIL174W} and \texttt{YLL017W}, are typed
\texttt{pseudogene} in R64 and fall outside our gene set because that set is
built from features of type \texttt{gene}. They are still real loci with SGD
records, \texttt{YLL017W} being SDC25, and Costanzo built and measured deletion
strains for both. Their absence is a property of how the gene set is
constructed, not of the genes.

\notesec{The same alias miss, a second time}
The build list's own training-support audit reports \texttt{YLR312C-B} as having
zero trigenic training records and labels it extrapolation. That count was taken
under the outdated name. The perturbed-gene index is keyed on current systematic
names, so the lookup returns nothing; under \texttt{YLR313C} the gene has 9
records. The gene on that panel with genuinely zero trigenic training records is
\texttt{YER079W}.
